% EE364B midterm progress report.
% Standard article style (11 pt, 1-in margins).
\documentclass[10pt,letterpaper]{article}

\usepackage[margin=1in]{geometry}
\usepackage{amsmath,amssymb,amsthm}
\usepackage{graphicx}
\usepackage{booktabs}
\usepackage{xcolor}
\usepackage[hidelinks]{hyperref}
\usepackage{cite}
\usepackage{parskip}

\newcommand{\Cref}[1]{Section~\ref{#1}}
\DeclareMathOperator*{\argmin}{arg\,min}
\newcommand{\R}{\mathbb{R}}
\newcommand{\opt}{^{\star}}

\begin{document}

\title{EE364B Progress Report:\\
Differentiable Inverse DC Optimal Power Flow}

\author{Carlo Schreiber}

\date{May 2026}

\maketitle

\section{Introduction}\label{sec:intro}

Grid operators publish hourly load and generation data, but rarely
release the associated cost curves and capacity limits that determine who gets
dispatched. I want to recover those parameters from public data. Knowing them enables counterfactual simulations, marginal
emission factor (MEF) estimates, and calibrated planning models
without access to confidential market filings.

To this extent, I frame the problem as inverse optimization~\cite{ahuja2001inverse,
esfahani2018data}: given dispatches that solve a DC optimal power
flow (DC-OPF), find the cost and capacity parameters that best
explain the data. The difficulty is that the forward problem is a
quadratic program (QP) with kinks in the parameters; whenever a
generator hits or leaves its limit, the gradient is undefined.

A standard approach~\cite{keshavarz2011imputing} sidesteps this by
fitting parameters to the KKT conditions of the observed dispatches.
It works for cost recovery, but breaks down for capacity recovery:
once the generator capacities $g_{\max}$ are free variables,
the complementarity conditions are trivially satisfied at any
differentiation~\cite{amos2017optnet}, which gives exact gradients
wherever the active set is locally constant. The approach handles
joint cost-and-capacity recovery and is available off-the-shelf
via \texttt{cvxpylayers}~\cite{agrawal2019differentiable}.

\section{Problem Formulation}\label{sec:problem}

Consider a network with $N$ buses, $L$ lines, and $G$ generators.
The forward DC-OPF I use is a QP in generator dispatches $g \in \R^G$
and line flows $p \in \R^L$:
\begin{equation}\label{eq:opf}
\begin{aligned}
\underset{g,p}{\text{minimize}}\quad
& f^{\!\top} g + \tfrac{\lambda}{2}\, g^{\!\top} g\\
\text{subject to}\quad
& A_g g + A_l p = d,\quad 0 \le g \le g_{\max},\quad |p| \le p_{\max},
\end{aligned}
\end{equation}
where $f \in \R^{G}_{+}$ is the vector of linear marginal costs,
$g_{\max}$ and $p_{\max}$ are generator and line capacity limits, and
$d \in \R^N$ is the demand. The small quadratic term
($\lambda = 10^{-3}$) keeps the dispatch unique when two generators
have the same cost.

Given $T$ noisy observations $\{(d_t, \hat g_t)\}_{t=1}^{T}$,
I want to recover the unknown parameters
$\theta\opt = (f\opt, g\opt_{\max})$ by solving
\begin{equation}\label{eq:inv}
\hat\theta = \argmin_{\theta \in \Theta}\;
\frac{1}{T}\sum_{t=1}^{T}
\bigl\| g\opt(d_t; \theta) - \hat g_t \bigr\|_2^{2}
+ R(\theta),
\end{equation}
where $g\opt(d_t; \theta)$ is the solution of \eqref{eq:opf} at demand
$d_t$, and $R(\theta)$ is a regulariser (Tikhonov plus an optional
$\ell_1$ merit-order term). I train with Adam and a cosine
learning-rate schedule over $300$--$600$ steps.

\paragraph{Identifiability.}
Not all parameters can be recovered from dispatch data alone.
The cost vector $f$ is only identifiable up to positive shifts within
the set of generators that are away from their bounds in at least one
sample; generators that are always at a limit never influence the
residual in \eqref{eq:inv}. Similarly, a capacity $g_{\max,i}$ is
only identifiable if generator $i$ saturates at least once.
I track both limits using a per-generator interior-fraction score
and examine the correlation with recovery error in \Cref{sec:results}.

\section{Methods}\label{sec:methods}

I implemented four estimators and compare them on the same
experiments:

\textbf{Ridge regression} and \textbf{MLP} (two hidden layers,
ReLU, width 128) both predict dispatch directly from demand.
They serve as forecasting baselines that completely ignore the
OPF structure and cannot recover $f$ or $g_{\max}$.

The \textbf{KKT-residual estimator}~\cite{keshavarz2011imputing,%
aswani2018inverse} fits $\theta$ by minimising the sum of squared
KKT residuals (stationarity and complementarity) evaluated at the
observed dispatches. This is a standard inverse-optimisation approach
and is my main point of comparison.

The \textbf{differentiable estimator} solves
\eqref{eq:inv} directly by backpropagating through \eqref{eq:opf}
via \texttt{cvxpylayers}. I run three variants:
\emph{caps fix} (capacities held at their true values, to isolate
cost recovery); \emph{full} (joint cost and capacity recovery);
and \emph{strat.} (full recovery with a Laplacian-smoothed
$24 \times G$ cost table to capture diurnal merit-order shifts).

For evaluation I use five metrics: clean-dispatch NRMSE on held-out
demands; cost cosine $\cos\hat f$ and Spearman rank correlation;
pairwise merit-order accuracy (fraction of generator pairs whose
cost ordering is correct); and capacity cosine $\cos\hat g_{\max}$.
Cosine and merit-order accuracy matter most since positive
scaling or shift of $f$ leaves the dispatch solution unchanged, so
there is a fundamental indeterminacy that RMSE on $f$ would not
reflect.

\section{Preliminary Results}\label{sec:results}

\subsection{Main Experiments}\label{sec:main}

A random 10-bus network ($G=10$, $L=12$, $T=200$, relative noise
$0.05$, three seeds) gives Table~\ref{tab:methods}.
The differentiable estimator with fixed capacities recovers cost
cosine $0.959$ versus $0.908$ for KKT ($+0.051$, $37\%$ lower
seed std). Joint capacity recovery costs one cosine point
($0.959 \to 0.949$) but recovers $\cos\hat g_{\max} \ge 0.999$
versus $0.870$ for KKT; this is the clearest evidence that differentiating
through the QP is worth the added complexity.

\begin{table}[h]
\centering
\caption{Headline Comparison, 10-Bus Synthetic Network.
  $T=200$, three seeds, mean$\,\pm\,$std.
  Ridge and MLP are forecasting baselines and do not estimate $f$.}
\label{tab:methods}
\setlength{\tabcolsep}{4pt}
\renewcommand{\arraystretch}{1.05}
\small
\begin{tabular}{@{}lcccc@{}}
\toprule
Method & NRMSE$_\text{cl.}$ & $\cos\hat f$ & Merit acc. & $\cos\hat g_{\max}$ \\
\midrule
Ridge              & $0.110$ & {--}            & {--}            & {--}            \\
MLP                & $0.477$ & {--}            & {--}            & {--}            \\
KKT                & $0.138$ & $.908{\pm}.019$ & $.881{\pm}.034$ & $.870{\pm}.055$ \\
Diff., caps fix.   & $\mathbf{0.125}$ & $\mathbf{.959{\pm}.012}$ & $\mathbf{.904{\pm}.013}$ & $\mathbf{1.00}$ \\
Diff., full        & $0.208$ & $.949{\pm}.018$ & $.889{\pm}.022$ & $.999$          \\
Diff., strat.      & $0.198$ & $.947{\pm}.018$ & $.874{\pm}.046$ & $.999$          \\
\bottomrule
\end{tabular}
\end{table}

For MEFs I computed $\partial(c^\top g\opt)/\partial d$ at $10$
operating points via autograd and compared against ground-truth
(Table~\ref{tab:mef}). Cosine $0.987$ is the strongest result in
the project; the gradient structure is preserved even where $\hat f$
is only recovered up to scale.

\begin{table}[h]
\centering
\caption{MEF Accuracy, 10 Operating Points.}
\label{tab:mef}
\setlength{\tabcolsep}{6pt}
\renewcommand{\arraystretch}{1.05}
\small
\begin{tabular}{@{}lcc@{}}
\toprule
Estimator & MEF cosine & MEF RMSE \\
\midrule
Diff., full      & $0.986$ & $0.070$ \\
Diff., caps fix. & $0.987$ & $0.064$ \\
\bottomrule
\end{tabular}
\end{table}

For diurnal data ($24$-hour merit-order rotation, renewables cheap
by day, peakers cheap at dusk) the static estimator recovers
$\cos\hat f = 0.963$ (NRMSE $0.365$) while the stratified estimator
gets $\cos\hat f = 0.956$ but also recovers the full hour-resolved
cost table at $\cos\hat F_\text{table} = 0.935$, which the static
model cannot represent.

On an 8-fuel single-bus PJM-like model ($198$~GW; load
$80$--$130$~GW; $T=400$; one seed) the estimator gives
$\cos\hat f = 0.672$, Spearman $0.810$, merit-order accuracy $0.786$
($22$ of $28$ fuel-pair orderings correct). The result rests on a
single seed; I will add two more before drawing conclusions.


\section{Next Steps}\label{sec:plan}

The main comparison (\Cref{sec:main}) and MEF result are in good
shape. The remaining work is about a few open
issues I want to resolve before submitting.

For the noise sweeps, my best guess is that the noise gets clipped
or rescaled somewhere in the data pipeline before it reaches training,
so the estimator never actually sees different noise levels. Once I
find and fix that I will rerun and check whether the cost cosine
drops meaningfully with $\sigma$. For the identifiability experiment,
I want to replace the interior-fraction proxy with a score based on
the rank of the active-set indicator matrix, which is more directly
connected to the identifiability argument. I also want to figure out
the finite-difference discrepancy; I plan to sweep the step size
$\varepsilon$ and rescale the parameters so the Jacobian entries are
comparable in magnitude.

Beyond those fixes, I would like two more seeds for PJM to get error
bars, a test on a 39-bus IEEE network, and a sweep of the regulariser
and Laplacian smoothing weights for the stratified estimator.

\bibliographystyle{plain}
\begin{thebibliography}{9}
\bibitem{ahuja2001inverse}
R.~K. Ahuja and J.~B. Orlin, ``Inverse Optimization,''
\emph{Operations Research}, vol.~49, no.~5, pp.~771--783, 2001.

\bibitem{keshavarz2011imputing}
A.~Keshavarz, Y.~Wang, and S.~Boyd, ``Imputing a Convex Objective
Function,'' in \emph{IEEE International Symposium on Intelligent
Control}, 2011, pp.~613--619.

\bibitem{esfahani2018data}
P.~M. Esfahani, S.~Shafieezadeh-Abadeh, G.~A. Hanasusanto, and
D.~Kuhn, ``Data-Driven Inverse Optimization with Imperfect
Information,'' \emph{Mathematical Programming}, vol.~167, no.~1,
pp.~191--234, 2018.

\bibitem{aswani2018inverse}
A.~Aswani, Z.-J. Shen, and A.~Siddiq, ``Inverse Optimization with
Noisy Data,'' \emph{Operations Research}, vol.~66, no.~3,
pp.~870--892, 2018.

\bibitem{amos2017optnet}
B.~Amos and J.~Z. Kolter, ``OptNet: Differentiable Optimization as a
Layer in Neural Networks,'' in \emph{International Conference on
Machine Learning}, 2017, pp.~136--145.

\bibitem{agrawal2019differentiable}
A.~Agrawal, B.~Amos, S.~Barratt, S.~Boyd, S.~Diamond, and J.~Z.
Kolter, ``Differentiable Convex Optimization Layers,'' in
\emph{Advances in Neural Information Processing Systems}, 2019,
pp.~9558--9570.

\bibitem{diamond2016cvxpy}
S.~Diamond and S.~Boyd, ``CVXPY: A Python-Embedded Modeling Language
for Convex Optimization,'' \emph{Journal of Machine Learning
Research}, vol.~17, no.~83, pp.~1--5, 2016.

\bibitem{boyd2004convex}
S.~Boyd and L.~Vandenberghe, \emph{Convex Optimization}.
Cambridge, U.K.: Cambridge University Press, 2004.
\end{thebibliography}

\end{document}
